
%%%%% 1. GESTIONE LINGUA, CARATTERI E CODIFICA
\usepackage[T1]{fontenc}
\usepackage[utf8]{inputenc}
\usepackage[english, italian]{babel}    % Gestisce entrambe le lingue, l'ultima è quella attiva di default
\usepackage{lmodern}                    % Carica i font stabili suggeriti dal template
\usepackage{hyphenat}                   % Gestione della sillabazione all'interno del testo (usare \nohypenation{parola} per evitare che venga spezzata)
\usepackage{csquotes}                   % Gestione del virgolettato all'interno delle citazioni (usare \enquote{testo} per inserire il testo tra virgolette)
\usepackage{mfirstuc}                   % Gestione della prima lettera maiuscola per parole selezionate (usare \makefirstuc{parola})

%%%%% 2. COLORI (Caricato subito con le opzioni del tuo file)
\usepackage[table, xcdraw]{xcolor}
\usepackage{colortbl}
\definecolor{blupolito}{RGB}{0, 46, 95}

%%%%% 3. GEOMETRIA E SPAZIATURA PAGINA
% Per forzare i margini con la dimensione voluta, togliere il simbolo "%" da \usepackage{geometry} e \geometry{...} qui sotto.
\usepackage{geometry}
%\geometry{a4paper, left = 20mm, top = 25mm, bottom = 25mm}
\usepackage{layout}                     % Utile per visualizzare graficamente il layout completo della pagina (usare \layout) --> solo diagnostica
\usepackage{setspace}                   % Gestione dell'interlinea all'interno del documento (usare \onehalfspacing o \setstretch{1.4})
\usepackage{ragged2e}                   % Miglioramento del controllo sull'allineamento del testo
\usepackage[skip=10pt minus 3pt, tocskip=5pt, indent=0pt]{parskip}     % Gestione della spaziatura tra i paragrafi (skip valore di spazio nominale, il minus indica che se serve per far stare il testo in pagina questo può essere ridotto di 3pt, tocskip è lo spazio nella tabella dei contenuti)

%%%%% 4. GRAFICA, IMMAGINI E CODICE
\usepackage{graphicx}                   % Gestione delle immagini in Latex
\usepackage[export]{adjustbox}          % Estensione funzionalità di centramento e ridimensionamento anche ad altra strutture come tabelle e formule
\usepackage{wrapfig}                    % Gestione della posizione fluida con testo dell'immagine (testo tutto intorno)
\usepackage{subfig}                     % Gestione dei grafici posti all'interno di un'unica figura
\usepackage{svg}                        % Gestione delle immagini in formato svg, garantendo una qualità perfetta
\usepackage{listings}                   % Gestione del codice sorgente, con possibilità di evidenziare la sintassi e di formattare il testo
\usepackage{verbatim}                   % Permette di inserire blocchi di test senza che vengano interpretati come codice Latex

%%%%% 5. MATEMATICA E FISICA
\usepackage{amsmath}                    % Pacchetto per la gestione avanzata della matematica in Latex
\usepackage{amsfonts}                   % Pacchetto per la gestione di font matematici aggiuntivi  
\usepackage{amssymb}                    % Pacchetto per la gestione di font matematici aggiuntivi
\usepackage{bm}                         % Pacchetto per la gestione di simboli matematici in grassetto           
\usepackage{mathtools}                  % Pacchetto per la gestione di simboli matematici (introduce ad esempio \xrightarrow{})
\usepackage{nicefrac}                   % Pacchetto per la gestione di frazioni oblique e compatte (usare \nicefrac{num}{den})
\usepackage{siunitx}                    % Pacchetto per la gestione delle unità di misura (usare \qty{20}{\kilogram\per\cubic\meter})

%%%%% 6. STRUTTURA DELLE TABELLE
\usepackage{array}                      % Pacchetto per la gestione avanzata delle tabelle in Latex
\usepackage{tabularx}                   % Pacchetto per la gestione di tabelle con colonne a larghezza fissa e altre a larghezza variabile
\usepackage{multirow}                   % Pacchetto che permette di unire celle verticalmente in una tabella
\usepackage{longtable}                  % Pacchetto per l'ambiente longtable per la creazione di tabelle che si estendono su più pagine
\usepackage{makecell}                   % Pacchetto per la gestione avanzata delle celle di una tabella, con possibilità di andare a capo
\usepackage{booktabs}                   % Pacchetto per la gestione delle linee orizzontali e verticali in una tabella (\toprule, \midrule, \bottomrule)
\usepackage{ltablex}                    % Pacchetto che unisce tabularx e longtable

%%%%% 7. FORMATTAZIONE CAPITOLI, TITOLI E DIDASCALIE
\usepackage{titlesec}                   % Pacchetto per la gestione avanzata dei titoli di capitolo, sezione e sottosezione
\usepackage[labelfont=bf]{caption}      % Pacchetto per la gestione avanzata delle didascalie di figure e tabelle
\titleformat{\chapter}{\normalfont\huge\bfseries\color{blupolito}}{\thechapter}{1em}{}
\titleformat{\section}{\normalfont\Large\bfseries\color{blupolito}}{\thesection}{1em}{}
\titleformat{\subsection}{\normalfont\large\bfseries\color{blupolito}}{\thesubsection}{1em}{}

%%%%% 8. GESTIONE TABELLE (Definizioni personalizzate dell'utente)
\newcolumntype{P}[1]{>{\centering\arraybackslash}m{#1}} 
\renewcommand{\cellalign}{cc} 
\newcolumntype{Q}[1]{>{\columncolor[HTML]{87cefa}\centering\arraybackslash}m{#1}} 
\renewcommand{\arraystretch}{0.5} 

%%%%% 9. ALGORITMI ED ELENCHI
\usepackage{algorithm}                  % Pacchetto per la gestione dell'ambiente algoritmo (introduce automaticamente anche l'indice apposito)
\usepackage{algpseudocode}              % Pacchetto che introduce un ambiente per scrivere pseudocodice in Latex
\usepackage{enumitem}                   % Pacchetto che permette il controllo di personalizzazione degli elenchi puntati e numerati
\usepackage{lipsum}                     % Pacchetto che introduce il comando \lipsum per generare testo fittizio

%%%%% 10. BIBLIOGRAFIA E CITAZIONI
\usepackage[final]{microtype}           % Rende la spaziatura del testo e della bibliografia perfetta
\usepackage[style=numeric, backend=biber, maxbibnames=9, maxcitenames=2, sorting=none, language=italian, autocite=superscript]{biblatex}           
\let\cite\autocite 

% Configurazione dell'ambiente bibliografia (dal tuo file)
\defbibenvironment{bibliography}
  {\list
     {\printfield[labelnumberwidth]{labelnumber}}
     {\setlength{\labelwidth}{\labelnumberwidth}%
      \setlength{\leftmargin}{\labelwidth}%
      \setlength{\labelsep}{\biblabelsep}%
      \addtolength{\leftmargin}{\labelsep}%
      \setlength{\itemsep}{\bibitemsep}%
      \setlength{\parsep}{\bibparsep}}%
      \renewcommand*{\makelabel}[1]{##1\hss}}
  {\endlist}
  {\item}     

%%%%% 11. COLLEGAMENTI IPERTESTUALI (Sempre per ultimo)
\usepackage{hyperref}                   % Pacchetto per la gestione dei collegamenti ipertestuali (link) all'interno del documento
\usepackage{bookmark}                   % Pacchetto per la gestione dei segnalibri (bookmarks) all'interno del documento
\usepackage{textcomp}                   % Pacchetto per l'introduzione di simboli extra (€, ©, ecc.) prima di hyperref
\usepackage{url}                        % Pacchetto per la gestione corretta dei link internet all'interno di testo e bibliografia

%%%%% 12. GLOSSARI E ACRONIMI
\usepackage[%
    toc,                                % Permette l'inserimento del glossario nell'indice 
    abbreviations,                      % Permette l'inserimento degli acronimi nel glossario
    nonumberlist,                       % Permette di non inserire i numeri di pagina delle occorrenze degli acronimi nel glossario
]{glossaries-extra}                     % Pacchetto per la gestione degli acronimi e dei glossari


%\usepackage{xcolor} % to define colors and use standard CSS names add dvipsnames as option, but it clashes with xcolor loaded in toptesi, pay attention that if it goes in conflict with tikz/beamer, simply use \documentclass[usenames,dvipsnames]{beamer}, along with other custom options when defining the document class
